% This file is intended to be included after amsmath, amssymb, and amsthm.
\section{Current-normalization gap}

Let $h:\mathbb{R}\to\mathbb{R}$ be an odd polynomial and let
$\mathcal H_h$ be its rectangular spectral map.  Equivalently, if
$h(s)=\sum_{k=0}^q a_ks^{2k+1}$, then
\begin{equation}
  \mathcal H_h(X)=\sum_{k=0}^q a_kX(X^\top X)^k.
\end{equation}
This is a smooth polynomial map even at repeated or zero singular values.  On
the Frobenius Hilbert space, define
\begin{equation}
  F_h(M)=\mathcal H_h\!\left(\frac{M}{\lVert M\rVert_F}\right),
  \qquad M\ne0.
\end{equation}

\begin{theorem}[Current-normalization gap]
  Fix $M\ne0$, write $r=\lVert M\rVert_F$ and $X=M/r$, and let
  $s_1,\ldots,s_d$ be the active singular values of $X$.  If
  $h'(s_i)\ne h'(s_j)$ for two active singular values, then
  $\operatorname{Sym}DF_h(M)$ is indefinite.  Consequently $F_h$ is not
  monotone on any open neighborhood of $M$.
\end{theorem}

\begin{proof}
  Let $A_X=D\mathcal H_h(X)$.  The polynomial spectral map is the gradient of
  \begin{equation}
    \Phi_h(X)=\sum_{k=0}^q
      \frac{a_k}{2k+2}\operatorname{tr}\!\left((X^\top X)^{k+1}\right),
  \end{equation}
  so $A_X$ is self-adjoint.  Let
  $Q_X[H]=H-\langle X,H\rangle_FX$ denote projection onto $X^\perp$.
  Differentiating current normalization gives
  \begin{equation}
    DF_h(M)=\frac{1}{r}A_XQ_X,
    \qquad
    S_M:=\operatorname{Sym}DF_h(M)
       =\frac{A_XQ_X+Q_XA_X}{2r}.
  \end{equation}
  If $X=U\operatorname{Diag}(s_i)V^\top$, radial differentiation yields
  \begin{equation}
    A_XX=U\operatorname{Diag}(s_ih'(s_i))V^\top.
  \end{equation}
  Thus the slope mismatch implies $w:=Q_XA_XX\ne0$.  Yet
  $\langle X,S_MX\rangle_F=0$ and $S_MX=w/(2r)\ne0$.  A symmetric positive-
  or negative-semidefinite map whose quadratic form vanishes at $X$ must
  annihilate $X$.  Hence $S_M$ is indefinite.  The differential necessary
  condition for monotonicity gives the conclusion.
\end{proof}

\begin{corollary}[Nonlinear polynomials in matrix rank at least three]
  If $d=\min\{m,n\}\ge3$ and $h$ is not affine, then $F_h$ is nonmonotone.
\end{corollary}

\begin{proof}
  Since $h'$ is a nonconstant polynomial, choose
  $a,b\in(0,1/\sqrt{2})$ with $h'(a)\ne h'(b)$ and set
  $c=\sqrt{1-a^2-b^2}>0$.  A matrix with normalized singular spectrum
  $(a,b,c,0,\ldots)$ satisfies the theorem.  In rank two the sharp slope
  condition remains valid, but exceptional nonlinear polynomials can have
  equal derivatives at every complementary pair, so the dimension qualifier
  cannot simply be dropped.
\end{proof}

On positive diagonal matrices the same result has an elementary coordinate
form.  Write $x=ru$, $D=\operatorname{Diag}(h'(u_i))$, and $P=I-uu^\top$.
Then $JF_h(x)=DP/r$ and $\operatorname{Sym}JF_h(x)=(DP+PD)/(2r)$.  For two
coordinates this is
\begin{equation}
  S(x)=\frac{1}{r}
  \begin{bmatrix}
    d_1u_2^2 & -\tfrac12(d_1+d_2)u_1u_2\\
    -\tfrac12(d_1+d_2)u_1u_2 & d_2u_1^2
  \end{bmatrix},
\end{equation}
and therefore
\begin{equation}
  \det S(x)
  =-\frac{u_1^2u_2^2}{4r^2}(d_1-d_2)^2<0.
\end{equation}

\begin{proposition}[Minimal constant linear shift]
  On any domain $\mathcal D$, define the pairwise deficit
  \begin{equation}
    \delta_{\mathrm{pair},\mathcal D}(T)
    =\sup_{M\ne N\in\mathcal D}
      \left[
      -\frac{\langle T(M)-T(N),M-N\rangle_F}
      {\lVert M-N\rVert_F^2}\right]_+.
  \end{equation}
  Then $T_\rho(M)=T(M)+\rho M$ is monotone on $\mathcal D$ if and only if
  $\rho\ge\delta_{\mathrm{pair},\mathcal D}(T)$.  If $T$ is continuously
  differentiable on an open convex $\mathcal D$, this pairwise deficit equals
  \begin{equation}
    \delta_{J,\mathcal D}(T)
    =\sup_{M\in\mathcal D}
      \max\!\left\{0,
      -\lambda_{\min}\!\left(\operatorname{Sym}JT(M)\right)\right\}.
  \end{equation}
\end{proposition}

\begin{proof}
  For each pair, the repair adds exactly $\rho\lVert M-N\rVert_F^2$ to the
  monotonicity product, proving the first statement.  The symmetric Jacobian
  of $T_\rho$ is $\operatorname{Sym}JT(M)+\rho I$.  On an open convex domain,
  integrating this Jacobian along line segments proves sufficiency of the
  local supremum, while directional limits of pairs prove necessity.
\end{proof}

\begin{proposition}[Scale obstruction]
  Suppose $T(\alpha M)=T(M)$ for every $\alpha>0$.  If $T$ has one violating
  pair, then its pairwise deficit on the nonzero ambient space is infinite.
  Hence no finite constant $\rho$ makes $T+\rho I$ globally monotone there.
\end{proposition}

\begin{proof}
  Scale both members of a violating pair by $\alpha$.  The numerator of the
  deficit ratio scales by $\alpha$, while its squared-distance denominator
  scales by $\alpha^2$.  The positive deficit therefore scales as
  $1/\alpha$ and diverges as $\alpha\downarrow0$.
\end{proof}
